% Auto-generated by results/main_results/scripts/generate_main_results.py
% Suggested packages: booktabs, threeparttable

\begin{table}[t]
\centering
\small
\begin{threeparttable}
\caption{Graph-granularity ablation on the two headline decoder settings.}
\label{tab:flashsvd-graph-ablation}
\begin{tabular}{lrrrr}
\toprule
Config & No graph & Split graph & Per-layer graph & No graph $\rightarrow$ per-layer \\
\midrule
512/32 & 32.424 & 16.587 & 9.913 & 3.27$\times$ \\
2048/128 & 29.938 & 16.401 & 10.102 & 2.96$\times$ \\
\bottomrule
\end{tabular}
\begin{tablenotes}[flushleft]\footnotesize
\item Values are decode latency in ms/token. Per-layer graph is the dominant systems improvement in the current runtime.
\end{tablenotes}
\end{threeparttable}
\end{table}

\begin{table}[t]
\centering
\small
\begin{threeparttable}
\caption{Runtime bookkeeping counts on the short-profile motivation study (prompt length 512, decode length 32).}
\label{tab:flashsvd-fragmentation}
\begin{tabular}{lrrrrr}
\toprule
Mode & Decode ms/token & Launches/token & Copies/token & Launch CPU ms/token & Copy CPU ms/token \\
\midrule
No graph & 36.004 & 1174 & 601 & 6.928 & 2.802 \\
Split graph & 20.533 & 694 & 793 & 4.647 & 5.565 \\
Per-layer & 9.355 & 86 & 249 & 2.076 & 2.708 \\
\bottomrule
\end{tabular}
\begin{tablenotes}[flushleft]\footnotesize
\item Launches aggregates direct CUDA kernel launches and CUDA graph launches. Copies aggregates copy, clone, and to-copy operators.
\end{tablenotes}
\end{threeparttable}
\end{table}

\begin{table}[t]
\centering
\small
\begin{threeparttable}
\caption{Attention-route microbenchmark for the current-token decode step.}
\label{tab:flashsvd-attn-route}
\begin{tabular}{rrrrr}
\toprule
Cached length & FlashSVD v1.5 & Dense + FA2 & Sparse + FA2 & Sparse legacy \\
\midrule
512 & 0.1016 & 0.2075 & 0.4389 & 4.3270 \\
2048 & 0.1183 & 0.2137 & 0.4509 & 7.9756 \\
4096 & 0.1390 & 0.2114 & 0.7313 & 13.4465 \\
\bottomrule
\end{tabular}
\begin{tablenotes}[flushleft]\footnotesize
\item All values are step latency in ms. Lower is better.
\end{tablenotes}
\end{threeparttable}
\end{table}

\begin{table}[t]
\centering
\small
\begin{threeparttable}
\caption{Kernel-level micro-ablations retained in the paper-facing summary.}
\label{tab:flashsvd-kernel-micro}
\begin{tabular}{llr}
\toprule
Component & Variant & Mean latency (ms) \\
\midrule
Attention reconstruct & Exact & 0.2522 \\
Attention reconstruct & Packed linear & 0.1484 \\
Attention reconstruct & Packed Triton & 0.1278 \\
\midrule
MLP backend & Baseline eager & 0.2030 \\
MLP backend & Prod graph & 0.2018 \\
\bottomrule
\end{tabular}
\begin{tablenotes}[flushleft]\footnotesize
\item The attention reconstruct path benefits materially from packed kernels; the MLP backend change is small in isolation.
\end{tablenotes}
\end{threeparttable}
\end{table}
